% Standalone problem document, generated from the adjacent README.md.
% Compile with XeLaTeX. Mathematical content is maintained in Markdown.
\documentclass[11pt,a4paper]{article}
\usepackage[fontsize=10.5pt]{scrextend}
\usepackage[margin=22mm,headheight=15pt,headsep=9mm,footskip=11mm]{geometry}
\usepackage{fontspec}
\setmainfont{texgyrepagella}[Extension=.otf,UprightFont=*-regular,BoldFont=*-bold,ItalicFont=*-italic,BoldItalicFont=*-bolditalic]
\setsansfont{texgyreheros}[Extension=.otf,UprightFont=*-regular,BoldFont=*-bold,ItalicFont=*-italic,BoldItalicFont=*-bolditalic]
\setmonofont{texgyrecursor}[Extension=.otf,UprightFont=*-regular,BoldFont=*-bold,ItalicFont=*-italic,BoldItalicFont=*-bolditalic,Scale=MatchLowercase]
\usepackage{amsmath,amssymb}
\usepackage{unicode-math}
\setmathfont{texgyrepagella-math.otf}
\usepackage{xcolor,microtype,enumitem,needspace,fancyhdr}
\usepackage{bookmark}
\definecolor{ink}{HTML}{17344B}
\definecolor{accent}{HTML}{197D80}
\definecolor{muted}{HTML}{596773}
\hypersetup{colorlinks=true,linkcolor=accent,urlcolor=accent,pdftitle={AC-09 - Deterministic
polynomial-time commutative Edmonds
problem},pdfauthor={Open Problems in Numerical Linear Algebra}}
\urlstyle{same}
\setlength{\parindent}{0pt}
\setlength{\parskip}{5pt plus 1pt minus 1pt}
\linespread{1.04}
\setlength{\emergencystretch}{3em}
\widowpenalty=10000
\clubpenalty=10000
\displaywidowpenalty=10000
\allowdisplaybreaks[1]
\setlist{nosep,leftmargin=1.5em}
\providecommand{\tightlist}{\setlength{\itemsep}{0pt}\setlength{\parskip}{0pt}}
\providecommand{\pandocbounded}[1]{#1}
\pagestyle{fancy}
\fancyhf{}
\fancyhead[L]{\sffamily\footnotesize\color{muted}OPEN PROBLEMS IN NUMERICAL LINEAR ALGEBRA}
\fancyhead[R]{\sffamily\bfseries\footnotesize\color{accent}AC-09}
\fancyfoot[L]{\parbox[b]{0.9\textwidth}{\sffamily\fontsize{8}{10}\selectfont\color{muted}Editorial ratings; a bounded search cannot certify that a problem remains unsolved.\\Literature check: 2026-09-10}}
\fancyfoot[R]{\sffamily\footnotesize\color{muted}\thepage}
\renewcommand{\headrulewidth}{0.3pt}
\renewcommand{\headrule}{\hbox to\headwidth{\color{accent}\leaders\hrule height \headrulewidth\hfill}}
\makeatletter
\renewcommand{\section}{\Needspace{5\baselineskip}\@startsection{section}{1}{0pt}{12pt plus 2pt minus 2pt}{4pt}{\sffamily\bfseries\large\color{ink}}}
\renewcommand{\subsection}{\Needspace{4\baselineskip}\@startsection{subsection}{2}{0pt}{9pt plus 1pt minus 1pt}{3pt}{\sffamily\bfseries\normalsize\color{ink}}}
\makeatother
\setcounter{secnumdepth}{0}
\begin{document}
{\sffamily\small\color{muted}Arithmetic and complexity\par}
\vspace{3pt}
{\raggedright\hyphenpenalty=10000\exhyphenpenalty=10000\sffamily\bfseries\fontsize{21}{24}\selectfont\color{ink}Deterministic
polynomial-time commutative Edmonds problem\par}
\vspace{7pt}
{\sffamily\small\textbf{Difficulty:} extreme\quad\textbf{Importance:} broadly
interesting\par}
{\sffamily\small\bfseries\color{accent}Status: Partially resolved\par}
\vspace{4pt}
{\color{accent}\hrule height 0.8pt}
\vspace{6pt}
\textbf{Rating rationale:} Extreme because deterministic symbolic
nonsingularity is a central derandomization barrier; broad importance
links matrix-space rank, polynomial identity testing and complexity
theory.\\
\textbf{Provenance:} source-stated problem, with rational input and bit
complexity made explicit.

The input consists of positive integers \(n,s\) and matrices
\(A_1,\ldots,A_s\in\mathbb Q^{n\times n}\), with rational entries
represented by signed binary numerators and positive binary
denominators. Does there exist a deterministic algorithm, running in a
number of bit operations bounded by a polynomial in the total binary
input length, that decides whether \[
\det\!\left(\sum_{i=1}^s x_iA_i\right)
\] is the zero polynomial in the commuting indeterminates
\(x_1,\ldots,x_s\) over \(\mathbb Q\)?

Equivalently, decide whether the rational linear space \[
\left\{\sum_{i=1}^s c_iA_i:\ c_1,\ldots,c_s\in\mathbb Q\right\}
\] contains a nonsingular matrix. No prescribed matrix structure is
assumed.

This is exact rank and nonsingularity testing for a parametrized matrix
space. Random evaluation gives a randomized route; the problem asks
whether randomness can be eliminated while retaining polynomial running
time.

\subsection{References}\label{references}

\begin{enumerate}
\def\labelenumi{\arabic{enumi}.}
\tightlist
\item
  C. Chindris and D. Kline, \emph{Edmonds' problem and the membership
  problem for orbit semigroups of quiver representations},
  arXiv:2008.13648 (2020), introduction: the general deterministic
  nonsingularity problem and the special classes treated in the paper.
  \href{https://arxiv.org/html/2008.13648v1}{Paper}.
\item
  G. Ivanyos, M. Karpinski, Y. Qiao, and M. Santha, \emph{Generalized
  Wong sequences and their applications to Edmonds' problems},
  arXiv:1307.6429v2 (2014), abstract and introduction: symbolic
  determinant identity testing, including the rational field and
  algorithms for restricted matrix spaces.
  \href{https://arxiv.org/abs/1307.6429}{Paper}.
\item
  A. Chatterjee, S. Ghosh, R. Gurjar, R. Raj, and T. Thierauf,
  \emph{Bipartite Matching is in NC}, ECCC TR26-100, revision 2 (2026),
  Section 1.1, ``Non-commutative rank'': explicitly distinguishes the
  still-open deterministic commutative rank problem.
  \href{https://eccc.weizmann.ac.il/report/2026/100/}{Report and version
  history}.
\end{enumerate}

\subsection{Status check --- 2026-09-10}\label{status-check-2026-09-10}

Rechecked
\href{https://eccc.weizmann.ac.il/report/2026/100/download/}{Chatterjee
et al., ECCC TR26-100, §1.1}, and searched for later commutative Edmonds
algorithms. The July 2026 revision explicitly states that no
deterministic polynomial-time algorithm is known generally.
\href{https://arxiv.org/abs/1307.6429}{Ivanyos et al.} gives algorithms
for restricted matrix spaces, substantive solved subclasses.
Deterministic noncommutative-rank algorithms and bipartite matching do
not settle unrestricted commuting symbolic determinants. No full
resolution was located.
\end{document}
